\documentclass[11pt,reqno]{amsart}
\usepackage[T1]{fontenc}
\usepackage[utf8]{inputenc}
\usepackage{lmodern}
\usepackage{amsmath,amssymb,amsthm,mathtools}
\usepackage[colorlinks=true,linkcolor=blue,citecolor=blue]{hyperref}
\usepackage[nameinlink,capitalize]{cleveref}
\newtheorem{theorem}{Theorem}[section]
\newtheorem{lemma}[theorem]{Lemma}
\newtheorem{proposition}[theorem]{Proposition}
\newtheorem{corollary}[theorem]{Corollary}
\theoremstyle{definition}
\newtheorem{assumption}[theorem]{Assumption}
\theoremstyle{remark}
\newtheorem{remark}[theorem]{Remark}
\newcommand{\C}{\mathbb C}
\newcommand{\R}{\mathbb R}
\newcommand{\T}{\mathbb T}
\newcommand{\Id}{\mathrm{Id}}
\newcommand{\norm}[1]{\left\lVert#1\right\rVert}
\newcommand{\abs}[1]{\left\lvert#1\right\rvert}
\title{Technical Appendix to Paper II:\ Full-Frequency Block Inversion and Uniform Moving-Family Cancellation}
\author{Qian Qi}
\date{August 29, 2026}
\begin{document}
\maketitle

\section{The scalar--quotient block decomposition}

Let
\[
 (\mathcal L_{a,z}h)(y)=
 \sum_{i=1}^4w_i(a)e^{-z\tau_i(a)}
 h(s_{i-1}(a)+w_i(a)y),
 \qquad z=\sigma+ib.
\]
Constants are invariant:
\[
 \mathcal L_{a,z}1=\lambda_a(z)1,
 \qquad
 \lambda_a(z)=\sum_iw_i(a)e^{-z\tau_i(a)}.
\]
Let $\ell(h)=\int_\T h$ and let $H=\ker\ell$.  Write
$h=c1+h_0$ with $h_0\in H$.  Relative to
$\C1\oplus H$, the operator has the form
\begin{equation}\label{eq:block-L}
 \mathcal L_{a,z}=
 \begin{pmatrix}
 \lambda_a(z)&\beta_{a,z}\\
 0&A_{a,z}
 \end{pmatrix},
\end{equation}
where $A_{a,z}$ represents the induced quotient operator and
$\beta_{a,z}h_0=\ell(\mathcal L_{a,z}h_0)$.

\begin{lemma}[Uniform quotient inverse]\label{lem:quotient-inverse}
For every $r\ge1$ and $\operatorname{Re}z\ge0$,
\[
 \norm{A_{a,z}}_{W^{r,1}/\C1\to W^{r,1}/\C1}\le\rho<1.
\]
Consequently
\[
 \norm{(I-A_{a,z})^{-1}}\le(1-\rho)^{-1},
 \qquad
 \norm{\beta_{a,z}}\le C_r
\]
uniformly in $a$ and $z$.
\end{lemma}

\begin{proof}
The quotient norm is equivalent to the sum of derivative $L^1$ norms.  For
$j\ge1$,
\[
 (\mathcal L_{a,z}h)^{(j)}=
 \sum_iw_i^{j+1}e^{-z\tau_i}h^{(j)}(s_{i-1}+w_i y).
\]
Changing variables gives
\[
 \norm{(\mathcal L_{a,z}h)^{(j)}}_1
 \le\rho^j\norm{h^{(j)}}_1
\]
because $\abs{e^{-z\tau_i}}\le1$.  The first derivative gives the common
factor $\rho$.  The Neumann series gives the inverse.  The functional
$\beta_{a,z}$ is bounded by the $W^{r,1}$ norm and the uniform branch and roof
bounds.
\end{proof}

From \eqref{eq:block-L},
\begin{equation}\label{eq:block-inverse}
 (I-\mathcal L_{a,z})^{-1}=
 \begin{pmatrix}
 (1-\lambda_a(z))^{-1}&
 (1-\lambda_a(z))^{-1}\beta_{a,z}(I-A_{a,z})^{-1}\\
 0&(I-A_{a,z})^{-1}
 \end{pmatrix}.
\end{equation}
Thus high frequency reduces to a scalar lower bound.

\section{Diophantine phase separation}

Assume branch weights satisfy $w_i(a)\ge w_*>0$ and roofs lie in
$[\tau_*,\tau^*]$.  Suppose
\begin{equation}\label{eq:dio}
 \inf_{k\in\mathbb Z^4}
 \max_i\abs{b\tau_i(a)-2\pi k_i}
 \ge c_0(1+\abs b)^{-\nu}
\end{equation}
uniformly in $a$ and $\abs b\ge1$.

\begin{lemma}[Weighted unit-circle convexity]\label{lem:circle}
Let $w_i\ge w_*$, $\sum_iw_i=1$, and $\theta_i\in\R$.  There are constants
$c,C>0$, depending only on $w_*$ and the number of phases, such that
\begin{equation}\label{eq:circle-convexity}
 \abs{1-\sum_iw_ie^{-i\theta_i}}
 \ge c\min\left\{1,
 \inf_{k\in\mathbb Z^4}\max_i\abs{\theta_i-2\pi k_i}^2\right\}.
\end{equation}
\end{lemma}

\begin{proof}
If the left side is bounded below by a fixed small constant there is nothing to
prove.  Otherwise the weighted average lies in a small neighbourhood of the
extreme point $1$ of the closed unit disk.  Uniform convexity of the Euclidean
norm implies every phase point $e^{-i\theta_i}$ lies in an arc of length
$C\sqrt\varepsilon$ around $1$, where
$\varepsilon=\abs{1-\sum_iw_ie^{-i\theta_i}}$.  Hence, for suitable integers
$k_i$,
\[
 \max_i\abs{\theta_i-2\pi k_i}\le C\sqrt\varepsilon.
\]
Rearranging gives \eqref{eq:circle-convexity}.  Equivalently, one may start
from the exact identity
\[
 1-\abs{\sum_iw_ie^{-i\theta_i}}^2
 =2\sum_{i<j}w_iw_j(1-\cos(\theta_i-\theta_j))
\]
and combine it with control of the common phase by the real part of the
weighted average.
\end{proof}

\begin{corollary}[Imaginary-axis gap]\label{cor:imaginary-gap}
Under \eqref{eq:dio},
\[
 \abs{1-\lambda_a(ib)}
 \ge c_1(1+\abs b)^{-2\nu},
 \qquad \abs b\ge1.
\]
\end{corollary}

\begin{proof}
Apply \cref{lem:circle} to $\theta_i=b\tau_i(a)$ and then
\eqref{eq:dio}.
\end{proof}

\section{Extension to a right half-strip}

Set
\[
 \delta_b=c_1(1+\abs b)^{-2\nu}.
\]

\begin{lemma}[Two-case strip gap]\label{lem:strip-gap}
There are $\sigma_0,c_2>0$ such that for
$0\le\sigma\le\sigma_0$ and $\abs b\ge1$,
\begin{equation}\label{eq:strip-gap}
 \abs{1-\lambda_a(\sigma+ib)}\ge c_2\delta_b
\end{equation}
uniformly in $a$.
\end{lemma}

\begin{proof}
There are constants $c_*,C^*>0$ such that, for
$0\le\sigma\le\sigma_0$,
\begin{align}
 1-\sum_iw_i(a)e^{-\sigma\tau_i(a)}&\ge c_*\sigma,
 \label{eq:radial-loss}\\
 \abs{\lambda_a(\sigma+ib)-\lambda_a(ib)}&\le C^*\sigma.
 \label{eq:sigma-cont}
\end{align}
Choose $\eta>0$ so small that $C^*\eta\le1/2$.  If
$\sigma\le\eta\delta_b$, then \eqref{eq:sigma-cont} and
\cref{cor:imaginary-gap} give
\[
 \abs{1-\lambda_a(\sigma+ib)}
 \ge\delta_b-C^*\sigma\ge\delta_b/2.
\]
If $\sigma>\eta\delta_b$, then
\[
 \abs{1-\lambda_a(\sigma+ib)}
 \ge1-\abs{\lambda_a(\sigma+ib)}
 \ge1-\sum_iw_ie^{-\sigma\tau_i}
 \ge c_*\eta\delta_b.
\]
Take $c_2=\min\{1/2,c_*\eta\}$.
\end{proof}

\begin{theorem}[Full strip resolvent]\label{thm:strip-resolvent}
For $0\le\operatorname{Re}z\le\sigma_0$ and
$\abs{\operatorname{Im}z}\ge1$,
\[
 \norm{(I-\mathcal L_{a,z})^{-1}}_{W^{r,1}\to W^{r,1}}
 \le C_r(1+\abs{\operatorname{Im}z})^{2\nu}.
\]
\end{theorem}

\begin{proof}
Insert \cref{lem:quotient-inverse,lem:strip-gap} into the explicit block
inverse \eqref{eq:block-inverse}.
\end{proof}

\section{Parameter derivatives}

Let $R_{a,z}=(I-\mathcal L_{a,z})^{-1}$.  Repeated differentiation gives
\begin{equation}\label{eq:resolvent-word}
 \partial_a^NR_{a,z}
 =\sum C_{k_1,\ldots,k_j}
 R_{a,z}(\partial_a^{k_1}\mathcal L_{a,z})R_{a,z}
 \cdots
 (\partial_a^{k_j}\mathcal L_{a,z})R_{a,z},
\end{equation}
where $k_1+\cdots+k_j=N$ and $1\le j\le N$.

\begin{lemma}[Twisted letter bound]\label{lem:twisted-letter}
For $r,k\ge0$,
\[
 \norm{\partial_a^k\mathcal L_{a,z}h}_{W^{r,1}}
 \le C_{r,k}(1+\abs z)^k\norm h_{W^{r+k,1}}.
\]
\end{lemma}

\begin{proof}
Each parameter derivative either hits a branch width, an inverse branch, or a
roof multiplier.  A roof derivative produces one factor of $z$; an inverse-
branch derivative produces one spatial derivative of $h$; widths are affine in
the explicit model and uniformly bounded in the general finite-order model.
Leibniz' rule gives the stated polynomial and regularity loss.
\end{proof}

\begin{theorem}[Derivative word bound]\label{thm:derivative-word}
For every $N$,
\[
 \norm{\partial_a^NR_{a,z}}_{W^{r+N,1}\to W^{r,1}}
 \le C_{r,N}(1+\abs{\operatorname{Im}z})^{2\nu(N+1)+N}
\]
on the high-frequency strip.
\end{theorem}

\begin{proof}
A word in \eqref{eq:resolvent-word} contains $j+1\le N+1$ resolvents and
letters of total derivative order $N$.  Apply
\cref{thm:strip-resolvent,lem:twisted-letter} and sum the finite set of ordered
compositions of $N$.
\end{proof}

\section{Uniform temporal shear for open billiards}

Let $\tau_N\circ h_1$ and $\tau_N\circ h_2$ be the return-roof sums on two
inverse branches of a symbolic return map.  Assume
\[
 \abs{D_v(\tau_N\circ h_1-\tau_N\circ h_2)}\ge\kappa>0
\]
on a cylinder $U$ and uniformly in the obstacle parameter.

\begin{lemma}[Oscillatory branch cancellation]\label{lem:oscillatory}
There are $c,C>0$ such that, for every sufficiently large $\abs b$, one can
choose a subcylinder of symbolic diameter $C/\abs b$ on which the two phases
differ by an angle in $[c,\pi-c]$.  On that subcylinder the sum of the two
normalized inverse-branch contributions loses a fixed fraction in modulus.
\end{lemma}

\begin{proof}
Along the direction $v$, the derivative of the phase difference
$b(\tau_N\circ h_1-\tau_N\circ h_2)$ has modulus at least $\abs b\kappa$.
The bounded second derivative makes this phase monotone on a cylinder of size
$C/\abs b$.  Choose a subinterval on which the phase difference is separated
from $2\pi\mathbb Z$.  Uniform lower and upper bounds on the two branch weights
then give strict triangle-inequality cancellation.
\end{proof}

\begin{proposition}[Uniform Dolgopyat block]\label{prop:Dolgopyat-block}
Assume uniform expansion, distortion, nonintegrability shear, and a mixing
finite-type symbolic base.  In the $b$-weighted Hölder norm there are
$C_0,c_0>0$ and a logarithmic iterate
$n_b\le C_0\log(2+\abs b)$ such that
\[
 \norm{\widehat{\mathscr L}_{a,ib}^{\,n_b}f}_b
 \le(1-c_0)\norm f_b
\]
for every normalized twisted operator in the compact moving family.
\end{proposition}

\begin{proof}
Use \cref{lem:oscillatory} on a fixed positive proportion of cylinders.
Outside those cylinders, the Lasota--Yorke estimate controls the Hölder part.
The symbolic mixing property returns every cone element to a cancellation
cylinder in $O(\log\abs b)$ steps.  Uniform expansion, distortion, and shear
margins make the cone invariance and lost fraction independent of $a$.  This is
the Dolgopyat cone construction with all constants taken on one compact
parameter family.
\end{proof}

\begin{corollary}[Uniform high-frequency inverse and derivatives]
The normalized open-billiard transfer operator satisfies
\[
 \norm{(I-\widehat{\mathscr L}_{a,ib})^{-1}}_{C^\alpha\to C^\alpha}
 \le C\abs b^\eta.
\]
If the symbolic roof and potential are analytic in $a$, every finite parameter
derivative of the inverse exists and has a polynomial high-frequency bound.
\end{corollary}

\begin{proof}
Sum the geometric series in blocks of length $n_b$.  Since
$n_b=O(\log\abs b)$, conversion between the weighted and ordinary Hölder norms
gives a polynomial loss.  Differentiate the inverse by
\eqref{eq:resolvent-word}; analytic potential derivatives have at most
polynomial $b$ growth.
\end{proof}

\section{The triangular lattice threshold}

Let the triangular lattice have covolume $\sqrt3/2$ and shortest vector length
one.

\begin{proposition}[All primitive directions]
For a primitive lattice vector $v$, adjacent lattice lines parallel to $v$ are
separated by
\[
 d(v)={\sqrt3/2\over\abs v}.
\]
Thus $d(v)\le\sqrt3/2$, with equality exactly in the shortest primitive
directions.  Disks of radius $r$ block every straight corridor if and only if
$r\ge\sqrt3/4$; strict finite horizon and nonoverlap hold for
$\sqrt3/4<r<1/2$.
\end{proposition}

\begin{proof}
The parallelogram spanned by $v$ and the shortest transverse lattice spacing
has area equal to the lattice covolume, proving the spacing formula.  A strip
between two adjacent lines survives exactly when its width exceeds $2r$.
Checking the maximal spacing therefore checks every primitive direction.  If
arbitrarily long free flights existed without an infinite corridor, translate
their initial points to a compact fundamental cell and take a convergent
subsequence of initial positions and directions; the limit is an infinite free
line, a contradiction.  The upper radius bound is nearest-neighbour
nonoverlap.
\end{proof}

\section{Conclusion}

The full-frequency theorem is the combination of a scalar Diophantine gap and
a uniformly contracting quotient.  The right-strip extension is a two-case
argument, and repeated parameter differentiation is controlled by explicit
resolvent words.  In the open-billiard branch, a concrete temporal shear enters
the uniform Dolgopyat cone construction on the fixed symbolic space.

\bibliographystyle{plain}
\bibliography{references}
\end{document}
